\lstset{
  basicstyle=\ttfamily\small,
  breaklines=true,
  breakatwhitespace=true,
  columns=fullflexible
}

\section{Prompt Templates}
\label{sec:appendix_prompts}

This appendix reproduces the complete prompt templates used for all three extraction tasks.
All prompts are injected into a single \texttt{user} message (no system prompt); the model's chat template adds the system/assistant framing.
Both versions provide a JSON-formatted type schema to the vLLM constrained-decoding engine; the schema is \emph{not} included in the natural-language prompt itself.
We use greedy decoding (\texttt{temperature=0}, \texttt{seed=42}) throughout.

\medskip\noindent
\textbf{Notation.}
Placeholders are shown in \texttt{\{braces\}}.
Text \textbf{added in v2} (relative to v1) is marked with $\triangleright$.
Text \textbf{removed in v2} is marked with \sout{strikethrough}.

\subsection{Named Entity Recognition (Few-NERD)}

\paragraph{v1 prompt (with character offsets).}

\begin{small}
\begin{lstlisting}
Extract all named entities from the following text.
Only extract entities of the following types:
- {type_1}: {definition_1}
- {type_2}: {definition_2}
  ...

Text: {input_text}

Output as JSON: [{"text": "...", "type": "...",
                  "start": ..., "end": ...}]
\end{lstlisting}
\end{small}

\noindent
\textbf{v1 constrained-decoding schema:}
\begin{small}
\begin{lstlisting}
{"type": "array",
 "items": {"type": "object",
   "properties": {
     "text":  {"type": "string"},
     "type":  {"type": "string",
               "enum": [<sorted_types>]},
     "start": {"type": "integer"},
     "end":   {"type": "integer"}},
   "required": ["text","type","start","end"]}}
\end{lstlisting}
\end{small}

\paragraph{v2 prompt (text+type only).}

\begin{small}
\begin{lstlisting}
Extract all named entities from the following text.
Only extract entities of the following types:
- {type_1}: {definition_1}
- {type_2}: {definition_2}
  ...

Text: {input_text}
\end{lstlisting}
\end{small}
\vspace{-0.5em}
\noindent$\triangleright$~\textbf{Added in v2:}
\begin{small}
\begin{lstlisting}
Only extract entities that clearly match one of the
given types. If no entities of these types exist in
the text, output an empty list [].
\end{lstlisting}
\end{small}
\vspace{-0.5em}
\begin{small}
\begin{lstlisting}
Output as JSON: [{"text": "...", "type": "..."}]
\end{lstlisting}
\end{small}

\noindent
\textbf{v2 constrained-decoding schema:}
\begin{small}
\begin{lstlisting}
{"type": "array",
 "items": {"type": "object",
   "properties": {
     "text": {"type": "string"},
     "type": {"type": "string",
              "enum": [<sorted_types>]}},
   "required": ["text","type"]}}
\end{lstlisting}
\end{small}

\paragraph{v1 $\to$ v2 differences (NER).}
\begin{enumerate}[leftmargin=*,itemsep=2pt]
\item \textbf{Offset fields removed.} The \texttt{start}/\texttt{end} integer fields are dropped from both the natural-language instruction and the JSON schema. In v1, 96.5\% of predicted offsets were hallucinated (\S\ref{sec:model}).
\item \textbf{Format constraints added.} Two sentences are appended: (a)~``Only extract entities that clearly match one of the given types'' suppresses spurious predictions; (b)~``If no entities \ldots output an empty list~\texttt{[]}'' provides an explicit empty-output path.
\end{enumerate}

\subsection{Event Detection (MAVEN)}

\paragraph{v1 prompt (with character offsets).}

\begin{small}
\begin{lstlisting}
Extract all event triggers from the following text.
A trigger is the word or phrase that most clearly
expresses the event occurrence. Only extract triggers
of the following event types:
- {type_1}: {definition_1}
- {type_2}: {definition_2}
  ...

Text: {input_text}

Output as JSON: [{"text": "...", "type": "...",
                  "start": ..., "end": ...}]
\end{lstlisting}
\end{small}

\paragraph{v2 prompt (text+type only).}

\begin{small}
\begin{lstlisting}
Extract all event triggers from the following text.
A trigger is the word or phrase that most clearly
expresses the event occurrence. Only extract triggers
of the following event types:
- {type_1}: {definition_1}
- {type_2}: {definition_2}
  ...

Text: {input_text}
\end{lstlisting}
\end{small}
\vspace{-0.5em}
\noindent$\triangleright$~\textbf{Added in v2:}
\begin{small}
\begin{lstlisting}
Only extract triggers that clearly match one of the
given event types. If no events of these types exist
in the text, output an empty list [].
\end{lstlisting}
\end{small}
\vspace{-0.5em}
\begin{small}
\begin{lstlisting}
Output as JSON: [{"text": "...", "type": "..."}]
\end{lstlisting}
\end{small}

\noindent The JSON schemas for ED are identical to the NER schemas above; only the natural-language framing differs (``event triggers'' instead of ``named entities'').

\paragraph{v1 $\to$ v2 differences (ED).}
Same two changes as NER: offset fields removed, format constraints added.

\subsection{Relation Extraction (FewRel)}

RE uses only a v2-style prompt (no offset fields are applicable to classification).

\paragraph{Prompt.}

\begin{small}
\begin{lstlisting}
Given the following sentence and two entities,
classify the relation between them.

Sentence: "{sentence}"
Head entity: "{head_entity}"
Tail entity: "{tail_entity}"

Candidate relations: {type_1}, {type_2}, ...

Output a single JSON object:
  {"relation": "<relation_type>"}
If none of the candidate relations apply, output:
  {"relation": "none"}
\end{lstlisting}
\end{small}

\noindent
\textbf{Constrained-decoding schema:}
\begin{small}
\begin{lstlisting}
{"type": "object",
 "properties": {
   "relation": {"type": "string",
     "enum": [<sorted_types>, "none"]}},
 "required": ["relation"]}
\end{lstlisting}
\end{small}

\subsection{Type Definitions}
\label{sec:appendix_type_defs}

All prompts include per-type natural-language definitions sourced from the respective dataset documentation:
\begin{itemize}[leftmargin=*,itemsep=2pt]
\item \textbf{Few-NERD:} 66 fine-grained entity types with definitions derived from the Few-NERD type ontology \citep{fewnerd-2021}.
\item \textbf{MAVEN:} 168 event types with FrameNet-based definitions \citep{maven-2020}.
\item \textbf{FewRel:} 80 relation types with Wikidata property descriptions \citep{fewrel-2018}.
\end{itemize}

\noindent Type definitions are injected as ``\texttt{- \{type\_name\}: \{definition\}}'' lines in alphabetical order. The number of types per prompt varies by experimental condition (count${}= 1$ to full schema; \S\ref{sec:exp_design}).
